% Source: https://artofproblemsolving.com/community/c3h1771914_mock_sprint_8
% Stable identity: aops-topic-1771914-q01
The slope does not depend on the size of the hexagon, so let its side be $2$. Its side on the $y$-axis runs from $(0,t)$ to $(0,t+2)$, and the hexagon lies to the right of the axis. Going round from $(0,t)$, the next vertex below is $(0,t)+2(\cos(-30^\circ),\sin(-30^\circ))=(\sqrt3,t-1)$, the lowest vertex; it lies on the $x$-axis, so $t=1$. The vertices are
\[
(0,1),\ (0,3),\ (\sqrt3,4),\ (2\sqrt3,3),\ (2\sqrt3,1),\ (\sqrt3,0),
\]
and the centre is $C=(\sqrt3,2)$ (the midpoint of $(0,1)$ and $(2\sqrt3,3)$).

\begin{asy}
usepackage("amsmath");
size(5cm);
defaultpen(fontsize(10pt));
real s3=sqrt(3);
pair V1=(0,1),V2=(0,3),V3=(s3,4),V4=(2*s3,3),V5=(2*s3,1),V6=(s3,0),O=(0,0),C=(s3,2);
draw((-0.6,0)--(4.4,0),Arrow(TeXHead));draw((0,-0.6)--(0,5),Arrow(TeXHead));
draw(V1--V2--V3--V4--V5--V6--cycle,linewidth(0.9));
draw(O--(2.2*s3,4.4),blue+linewidth(0.9));
dot(C,blue);dot(V1);dot(V2);dot(V6);
label("$C$",C,dir(-20),blue);
label("$(0,1)$",V1,W);label("$(0,3)$",V2,W);label("$(\sqrt3,0)$",V6,SE);
label("$O$",O,SW);label("$x$",(4.4,0),E);label("$y$",(0,5),N);
label("\mbox{}",(min(currentpicture,user=true)+max(currentpicture,user=true))/2,black); // keeps the text after the figure black when LaTeX typesets the labels inline
\end{asy}

A regular hexagon is symmetric under the half-turn about its centre, so every line through $C$ cuts it into two congruent halves of equal area. A line through the origin not passing through $C$ is a translate of the line through $C$ parallel to it, moved off the centre, so one side gains area and the halves are unequal. Hence the required line is $OC$, with slope
\[
\dfrac{2}{\sqrt3}=\boxed{\dfrac{2\sqrt3}{3}} .
\]
